\chapter{Enums and Pattern Matching}
\label{ch:enums-matching}

Many values in a program are not really numbers or text they are a choice from
a small, fixed set of possibilities. A day is one of seven; a traffic light is
red, amber, or green; an order is pending, shipped, or delivered. An
\textbf{enum} (short for \emph{enumeration}) lets you name such a set and use those
names directly, instead of remembering that ``3 means Thursday.''

\section{Defining an enum}

An enum is declared with the \code{enum} keyword, a name, and a list of
\textbf{members} the allowed values closed with \code{end}. The members may
be written on one line, separated by commas, or one per line:

\begin{salam}
enum Day: Mon, Tue, Wed, Thu, Fri, Sat, Sun end
\end{salam}

This defines a new type, \code{Day}, whose only values are the seven named days.
A variable of type \code{Day} can hold exactly one of them nothing else. That
is the whole point: the type makes invalid values \emph{unrepresentable}.

\section{Using enum values}

You refer to a member by qualifying it with the enum's name and a dot:
\code{Day.Mon}, \code{Day.Sat}. This makes it clear which enum the value belongs
to and avoids name clashes between enums that happen to share a member name.

\begin{salam}
enum Day: Mon, Tue, Wed, Thu, Fri, Sat, Sun end

func main:
    today Day = Day.Sat
    if today == Day.Sat || today == Day.Sun:
        println "It's the weekend!"
    else:
        println "A working day."
    end
end
\end{salam}

\begin{progout}
It's the weekend!
\end{progout}

You compare enum values with \code{==} and \code{!=}, exactly as you would
numbers. The variable \code{today} can only ever be one of the seven days, so the
compiler and the next person to read your code knows precisely what values
are in play.

\section{Enums are numbers underneath}

Each enum member has an integer value. By default they count up from zero in
order: in \code{Day}, \code{Mon} is 0, \code{Tue} is 1, and so on to \code{Sun}
at 6. You can see a member's number by converting it to an integer with \code{as}:

\begin{salam}
enum Day: Mon, Tue, Wed, Thu, Fri, Sat, Sun end

func main:
    println Day.Mon as i32      // 0
    println Day.Sat as i32      // 5
end
\end{salam}

\begin{progout}
0
5
\end{progout}

\subsection{Choosing the numbers}

When the underlying numbers matter because they correspond to codes, flags, or
external values you can assign them explicitly. Any member you do not give a
number to continues counting up from the previous one:

\begin{salam}
enum Color:
    Red            // 0
    Green = 5      // 5
    Blue           // 6 (continues from 5)
end

func main:
    println Color.Red as i32, Color.Green as i32, Color.Blue as i32
end
\end{salam}

\begin{progout}
0 5 6
\end{progout}

\code{Red} keeps the default 0, \code{Green} is pinned to 5, and \code{Blue}
follows on at 6. This is handy when an enum mirrors values defined elsewhere, such
as HTTP status codes or hardware register values.

\begin{warn}
The conversion goes one way freely: an enum value turns into its integer with
\code{as}. Going the other way treating an arbitrary integer as an enum is
something to do only with care, because a number might not correspond to any
member. Prefer to keep values in their enum form and convert to integers only at
the edges of your program, where you must.
\end{warn}

\section{Pattern matching with \texttt{if}/\texttt{else if}}

Some languages provide a dedicated \code{match} or \code{switch} statement for
choosing an action based on which enum value you have. Salam keeps things simple:
you ``match'' on an enum with an ordinary \code{if}/\code{else if} chain, which
you already know well from Chapter~4.

\begin{salam}
enum Light: Red, Amber, Green end

func action(light Light) str:
    if light == Light.Red:
        ret "Stop"
    else if light == Light.Amber:
        ret "Slow down"
    else:
        ret "Go"
    end
end

func main:
    println action(Light.Red)
    println action(Light.Amber)
    println action(Light.Green)
end
\end{salam}

\begin{progout}
Stop
Slow down
Go
\end{progout}

The function takes a \code{Light} and returns the right instruction for it. The
\code{else if} chain considers each possibility in turn; the final \code{else}
catches the last case (\code{Green}). This pattern one branch per enum
member is the Salam way of ``matching'' on an enum.

\begin{tip}
When you write such a chain, cover \emph{every} member, and let the final
\code{else} handle the last one (or an unexpected value) rather than leaving a
case out. If you later add a member to the enum, revisit these chains to handle
it a forgotten case is a classic source of bugs in any language.
\end{tip}

\section{Enums as state machines}

Enums shine when a value moves through a sequence of well-defined states. Consider
an order that progresses from \code{Pending} to \code{Shipped} to
\code{Delivered}. An enum names the states, and a function decides the next one:

\begin{salam}
enum Status: Pending, Shipped, Delivered end

func next(s Status) Status:
    if s == Status.Pending:
        ret Status.Shipped
    else if s == Status.Shipped:
        ret Status.Delivered
    else:
        ret Status.Delivered      // already final; stays put
    end
end

func label(s Status) str:
    if s == Status.Pending:   ret "pending"
    else if s == Status.Shipped: ret "shipped"
    else:                     ret "delivered"
    end
end

func main:
    mut s Status = Status.Pending
    println label(s)
    s = next(s)
    println label(s)
    s = next(s)
    println label(s)
end
\end{salam}

\begin{progout}
pending
shipped
delivered
\end{progout}

The enum guarantees that \code{s} is always a valid status, and the \code{next}
function encodes the legal transitions in one place. This is a small example of a
\textbf{state machine} a powerful and common way to model anything that moves
through stages, from a game's screens to a network connection's lifecycle.

\begin{deep}[In Depth: why an enum beats a bare integer]
You could represent a status with a plain \code{i32} 0 for pending, 1 for
shipped, and so on and many programs do. But then nothing stops you from
assigning \code{99}, or from adding a \code{Status} to a \code{Day} by accident,
or from forgetting which number meant what. An enum makes those mistakes
impossible: the type system knows the only valid values, the names document
themselves, and the compiler checks that you use them consistently. The cost is
nothing an enum compiles down to exactly the integer you would have written by
hand so you get the safety for free.
\end{deep}

\section{Chapter summary}

\begin{itemize}[leftmargin=1.4em]
  \item An \code{enum} names a fixed set of values; a variable of that type holds
        exactly one member.
  \item Refer to members as \code{EnumName.Member}; compare them with \code{==}
        and \code{!=}.
  \item Members have integer values, counting from 0 by default; convert with
        \code{as i32}. You may assign explicit values, and later members continue
        counting up.
  \item Salam has no \code{match}/\code{switch}; ``match'' on an enum with an
        \code{if}/\code{else if} chain that covers every member.
  \item Enums are ideal for modelling states and transitions (state machines),
        and are safer than bare integers at no run-time cost.
\end{itemize}

\section{Exercises}

\exercise Define \code{enum Suit: Hearts, Diamonds, Clubs, Spades end} and print
the integer value of each member.

\exercise Write a function \code{is\_red(s Suit) bool} that returns \code{true} for
\code{Hearts} and \code{Diamonds}, and test it on all four suits.

\exercise Define \code{enum Direction: North, East, South, West end} and a function
\code{turn\_right(d Direction) Direction} that rotates clockwise (North becomes
East, and so on, West wrapping back to North).

\exercise Model a traffic light that cycles \code{Green} $\to$ \code{Amber} $\to$
\code{Red} $\to$ \code{Green}. Write a \code{next} function and run it through two
full cycles.

\exercise Define \code{enum HttpStatus} with \code{Ok = 200},
\code{NotFound = 404}, and \code{ServerError = 500}, and print each member's
number.
